% Options for packages loaded elsewhere
\PassOptionsToPackage{unicode}{hyperref}
\PassOptionsToPackage{hyphens}{url}
%
\documentclass[
]{article}
\usepackage{amsmath,amssymb}
\usepackage{iftex}
\ifPDFTeX
  \usepackage[T1]{fontenc}
  \usepackage[utf8]{inputenc}
  \usepackage{textcomp} % provide euro and other symbols
\else % if luatex or xetex
  \usepackage{unicode-math} % this also loads fontspec
  \defaultfontfeatures{Scale=MatchLowercase}
  \defaultfontfeatures[\rmfamily]{Ligatures=TeX,Scale=1}
\fi
\usepackage{lmodern}
\ifPDFTeX\else
  % xetex/luatex font selection
\fi
% Use upquote if available, for straight quotes in verbatim environments
\IfFileExists{upquote.sty}{\usepackage{upquote}}{}
\IfFileExists{microtype.sty}{% use microtype if available
  \usepackage[]{microtype}
  \UseMicrotypeSet[protrusion]{basicmath} % disable protrusion for tt fonts
}{}
\makeatletter
\@ifundefined{KOMAClassName}{% if non-KOMA class
  \IfFileExists{parskip.sty}{%
    \usepackage{parskip}
  }{% else
    \setlength{\parindent}{0pt}
    \setlength{\parskip}{6pt plus 2pt minus 1pt}}
}{% if KOMA class
  \KOMAoptions{parskip=half}}
\makeatother
\usepackage{xcolor}
\usepackage[margin=1in]{geometry}
\usepackage{color}
\usepackage{fancyvrb}
\newcommand{\VerbBar}{|}
\newcommand{\VERB}{\Verb[commandchars=\\\{\}]}
\DefineVerbatimEnvironment{Highlighting}{Verbatim}{commandchars=\\\{\}}
% Add ',fontsize=\small' for more characters per line
\usepackage{framed}
\definecolor{shadecolor}{RGB}{248,248,248}
\newenvironment{Shaded}{\begin{snugshade}}{\end{snugshade}}
\newcommand{\AlertTok}[1]{\textcolor[rgb]{0.94,0.16,0.16}{#1}}
\newcommand{\AnnotationTok}[1]{\textcolor[rgb]{0.56,0.35,0.01}{\textbf{\textit{#1}}}}
\newcommand{\AttributeTok}[1]{\textcolor[rgb]{0.13,0.29,0.53}{#1}}
\newcommand{\BaseNTok}[1]{\textcolor[rgb]{0.00,0.00,0.81}{#1}}
\newcommand{\BuiltInTok}[1]{#1}
\newcommand{\CharTok}[1]{\textcolor[rgb]{0.31,0.60,0.02}{#1}}
\newcommand{\CommentTok}[1]{\textcolor[rgb]{0.56,0.35,0.01}{\textit{#1}}}
\newcommand{\CommentVarTok}[1]{\textcolor[rgb]{0.56,0.35,0.01}{\textbf{\textit{#1}}}}
\newcommand{\ConstantTok}[1]{\textcolor[rgb]{0.56,0.35,0.01}{#1}}
\newcommand{\ControlFlowTok}[1]{\textcolor[rgb]{0.13,0.29,0.53}{\textbf{#1}}}
\newcommand{\DataTypeTok}[1]{\textcolor[rgb]{0.13,0.29,0.53}{#1}}
\newcommand{\DecValTok}[1]{\textcolor[rgb]{0.00,0.00,0.81}{#1}}
\newcommand{\DocumentationTok}[1]{\textcolor[rgb]{0.56,0.35,0.01}{\textbf{\textit{#1}}}}
\newcommand{\ErrorTok}[1]{\textcolor[rgb]{0.64,0.00,0.00}{\textbf{#1}}}
\newcommand{\ExtensionTok}[1]{#1}
\newcommand{\FloatTok}[1]{\textcolor[rgb]{0.00,0.00,0.81}{#1}}
\newcommand{\FunctionTok}[1]{\textcolor[rgb]{0.13,0.29,0.53}{\textbf{#1}}}
\newcommand{\ImportTok}[1]{#1}
\newcommand{\InformationTok}[1]{\textcolor[rgb]{0.56,0.35,0.01}{\textbf{\textit{#1}}}}
\newcommand{\KeywordTok}[1]{\textcolor[rgb]{0.13,0.29,0.53}{\textbf{#1}}}
\newcommand{\NormalTok}[1]{#1}
\newcommand{\OperatorTok}[1]{\textcolor[rgb]{0.81,0.36,0.00}{\textbf{#1}}}
\newcommand{\OtherTok}[1]{\textcolor[rgb]{0.56,0.35,0.01}{#1}}
\newcommand{\PreprocessorTok}[1]{\textcolor[rgb]{0.56,0.35,0.01}{\textit{#1}}}
\newcommand{\RegionMarkerTok}[1]{#1}
\newcommand{\SpecialCharTok}[1]{\textcolor[rgb]{0.81,0.36,0.00}{\textbf{#1}}}
\newcommand{\SpecialStringTok}[1]{\textcolor[rgb]{0.31,0.60,0.02}{#1}}
\newcommand{\StringTok}[1]{\textcolor[rgb]{0.31,0.60,0.02}{#1}}
\newcommand{\VariableTok}[1]{\textcolor[rgb]{0.00,0.00,0.00}{#1}}
\newcommand{\VerbatimStringTok}[1]{\textcolor[rgb]{0.31,0.60,0.02}{#1}}
\newcommand{\WarningTok}[1]{\textcolor[rgb]{0.56,0.35,0.01}{\textbf{\textit{#1}}}}
\usepackage{graphicx}
\makeatletter
\def\maxwidth{\ifdim\Gin@nat@width>\linewidth\linewidth\else\Gin@nat@width\fi}
\def\maxheight{\ifdim\Gin@nat@height>\textheight\textheight\else\Gin@nat@height\fi}
\makeatother
% Scale images if necessary, so that they will not overflow the page
% margins by default, and it is still possible to overwrite the defaults
% using explicit options in \includegraphics[width, height, ...]{}
\setkeys{Gin}{width=\maxwidth,height=\maxheight,keepaspectratio}
% Set default figure placement to htbp
\makeatletter
\def\fps@figure{htbp}
\makeatother
\setlength{\emergencystretch}{3em} % prevent overfull lines
\providecommand{\tightlist}{%
  \setlength{\itemsep}{0pt}\setlength{\parskip}{0pt}}
\setcounter{secnumdepth}{-\maxdimen} % remove section numbering
\usepackage{amsmath}
\usepackage{mathtools}
\usepackage{amsthm}
\usepackage{bm}
\usepackage{xcolor}
\usepackage{fbox}
\usepackage{amsmath}
\ifLuaTeX
  \usepackage{selnolig}  % disable illegal ligatures
\fi
\IfFileExists{bookmark.sty}{\usepackage{bookmark}}{\usepackage{hyperref}}
\IfFileExists{xurl.sty}{\usepackage{xurl}}{} % add URL line breaks if available
\urlstyle{same}
\hypersetup{
  pdftitle={Computational Statistics - STAT 575 - Exam \#1},
  pdfauthor={Alex Towell (atowell@siue.edu)},
  hidelinks,
  pdfcreator={LaTeX via pandoc}}

\title{Computational Statistics - STAT 575 - Exam \#1}
\author{Alex Towell
(\href{mailto:atowell@siue.edu}{\nolinkurl{atowell@siue.edu}})}
\date{}

\begin{document}
\maketitle

\newcommand{\var}{\operatorname{Var}}
\newcommand{\expect}{\operatorname{E}}
\newcommand{\cov}{\operatorname{Cov}}
\newcommand{\se}{\operatorname{SE}}
\newcommand{\mat}[1]{\mathbf{#1}} 
\newcommand{\eval}[2]{\left. #1 \right\vert_{#2}}
\newcommand{\poi}{\operatorname{POI}}
\newcommand{\argmax}{\operatorname{arg\,max}}

\hypertarget{problem-1}{%
\section{Problem 1}\label{problem-1}}

\hypertarget{part-1-10-points}{%
\subsection{Part (1) (10 points)}\label{part-1-10-points}}

Use accept-reject algorithm to generate a truncate standard Normal
distribution, with density \[
  f(x) ∝ e^{-x^2/2}I(x > 2).
\] Consider the proposed distribution \(g(x) = 2 e^{−2(x−2)}I(x > 2)\).

\hypertarget{part-a}{%
\subsection{Part (a)}\label{part-a}}

\fbox{Sample from $g$ using inverse-transform method.}

First, we must find the cdf \(G\) of \(g\). One way to derive \(G\) is
to notice that \(g\) is the density of the shifted exponential, i.e., if
\(S \sim \exp(\lambda=-2)\), then \(X = 2 + S\) has a cdf \begin{align}
  G(x) &= P(X \leq x)\\
       &= P(2+S \leq x)\\
       &= P(S \leq x-2)\\
       &= F_S(x-2)\\
       &= &= I(x > 2)(1 - \exp(-2(x-2))).
\end{align} As a quick proof that \(G\) is the cdf of \(g\), note that
\(dG/dx = g\).

Alternatively, the cdf \(G\) is defined as \begin{align}
  G(x) &= \int_{-\infty}^{x} g(s) ds\\
       &= I(x > 2) \int_{2}^{x} 2 e^{−2(s−2)} ds\\
       &= I(x > 2) e^4 \int_{2}^{x} 2 e^{−2 s} ds\\
       &= -I(x > 2) e^4 \int_{2}^{x} (-2) e^{−2 s} ds\\
       &= -I(x > 2) e^4 \left(\left. e^{-2 s} \right\vert_{2}^{x}\right)\\
       &= -I(x > 2) e^4 \left(e^{-2 x}-e^{-4}\right)\\
       &= I(x > 2) (1 - e^4 e^{-2 x})\\
       &= I(x > 2) (1 - e^{-2(x-2)}),
\end{align} which obtains the same result. Either way, we solve for
\(x\) in \begin{align*}
  u           &= G(x)\\
  u           &= 1 - e^{-2(x-2)}\\
  1-u         &= e^{-2(x-2)}\\
  \log(1-u)   &= -2(x-2)\\
  -\frac{1}{2}\log(1-u) &= x-2,
\end{align*} and thus \[
  x = -\frac{1}{2}\log(1-u) + 2
\] where \(u\) is drawn from \(\rm{UNIF}(0,1)\). That is, if
\(U \sim \rm{UNIF}(0,1)\) then \[
  X = -\frac{1}{2}\log(1-U) + 2
\] is a random variable whose density is \(g\).

\hypertarget{part-b}{%
\subsection{Part (b)}\label{part-b}}

\fbox{%
\begin{minipage}{.8\textwidth}
Use accept-reject algorithm to generate a sample of $10000$ from f. Choose a $c$
so that $f(y)/[c g(y)] ≤ 1$. Verify the generated sample via a plot of the true
normalized density, and a histogram of the generated values.
\end{minipage}}

We wish to find a \(c\) such that \(f(y)/[c g(y)] ≤ 1\).

However, we do not need \(f\), we only need the kernel of \(f\), which
has already been given.

So, ideally, we find the smallest \(c\) that satisifes the inequality \[
  \frac{k(y)}{g(y)} \leq 1
\] where \(k\) is the kernel of \(f\), which is given by
\(c = \max\{k(y)/g(y)\}\). We may solve this analytically by finding \[
  y^* = \operatorname{arg\,max}_y k(y)/g(y)
\] and then let \(c = k(y^*)/g(y^*)\).

First, we let \[
  h(y) = k(y)/g(y) = \frac{e^{-y^2/2}}{2e^{-2(y-2)}}.
\] We can find the value that maximizes \(h\) by finding the value that
maximizes \(\log h\), which is given by \[
  \log h(y) = \log \frac{e^{-y^2/2}}{2e^{-2(y-2)}} = -\frac{1}{2}y^2 + 2(y-2).
\] The value \(y^*\) that maximizes \(\log h\) satisfies \[
  \left. \frac{d \log h}{d y} \right\vert_{y=y^*} = -y^* + 2 = 0,
\] which means \(y^*=2\) maximizes \(\log h\) and thus \begin{align}
  c = h(2) &= \frac{k(2)}{g(2)}\\
           &= \frac{1}{2e^2}.
\end{align}

Thus, we sample \(y\) from \(g\) and \(u\) from \(\rm{UNIF}(0,1)\) and
accept the sample as being drawn from \(f\) if \[
  u \leq \frac{k(y)}{g(y)} = e^{-(y-2)^2/2}.
\]

Here is the code that implements our sampling method for \(f\) using
acceptance-rejection sampling:

\begin{Shaded}
\begin{Highlighting}[]
\CommentTok{\# the proposed density we can easily sample from}
\NormalTok{g }\OtherTok{\textless{}{-}} \ControlFlowTok{function}\NormalTok{(y)}
\NormalTok{\{}
  \DecValTok{2}\SpecialCharTok{*}\FunctionTok{exp}\NormalTok{(}\SpecialCharTok{{-}}\DecValTok{2}\SpecialCharTok{*}\NormalTok{(y}\DecValTok{{-}2}\NormalTok{))}
\NormalTok{\}}

\CommentTok{\# inverse transform method for sampling from g}
\NormalTok{rg }\OtherTok{\textless{}{-}} \ControlFlowTok{function}\NormalTok{(n)}
\NormalTok{\{}
  \SpecialCharTok{{-}}\FloatTok{0.5}\SpecialCharTok{*}\FunctionTok{log}\NormalTok{(}\DecValTok{1}\SpecialCharTok{{-}}\FunctionTok{runif}\NormalTok{(n))}\SpecialCharTok{+}\DecValTok{2}
\NormalTok{\}}

\NormalTok{rf }\OtherTok{\textless{}{-}} \ControlFlowTok{function}\NormalTok{(n)}
\NormalTok{\{}
\NormalTok{  data }\OtherTok{\textless{}{-}} \FunctionTok{vector}\NormalTok{(}\AttributeTok{length=}\NormalTok{n)}
  \ControlFlowTok{for}\NormalTok{ (i }\ControlFlowTok{in} \DecValTok{1}\SpecialCharTok{:}\NormalTok{n)}
\NormalTok{  \{}
    \ControlFlowTok{repeat}
\NormalTok{    \{}
\NormalTok{      y }\OtherTok{\textless{}{-}} \FunctionTok{rg}\NormalTok{(}\DecValTok{1}\NormalTok{)}
\NormalTok{      u }\OtherTok{\textless{}{-}} \FunctionTok{runif}\NormalTok{(}\DecValTok{1}\NormalTok{)}
      \ControlFlowTok{if}\NormalTok{ (}\FunctionTok{runif}\NormalTok{(}\DecValTok{1}\NormalTok{) }\SpecialCharTok{\textless{}=} \FunctionTok{exp}\NormalTok{(}\SpecialCharTok{{-}}\NormalTok{(y}\DecValTok{{-}2}\NormalTok{)}\SpecialCharTok{\^{}}\DecValTok{2}\SpecialCharTok{/}\DecValTok{2}\NormalTok{))}
\NormalTok{      \{}
\NormalTok{        data[i] }\OtherTok{\textless{}{-}}\NormalTok{ y}
        \ControlFlowTok{break}
\NormalTok{      \}}
\NormalTok{    \}}
\NormalTok{  \}}
\NormalTok{  data}
\NormalTok{\}}

\NormalTok{k }\OtherTok{\textless{}{-}} \ControlFlowTok{function}\NormalTok{(y)}
\NormalTok{\{}
  \FunctionTok{exp}\NormalTok{(}\SpecialCharTok{{-}}\NormalTok{y}\SpecialCharTok{\^{}}\DecValTok{2}\SpecialCharTok{/}\DecValTok{2}\NormalTok{)}
\NormalTok{\}}

\NormalTok{f }\OtherTok{\textless{}{-}} \ControlFlowTok{function}\NormalTok{(y)}
\NormalTok{\{}
  \CommentTok{\# normalizing constant}
\NormalTok{  Z }\OtherTok{\textless{}{-}} \FloatTok{17.5358}
\NormalTok{  Z}\SpecialCharTok{*}\FunctionTok{k}\NormalTok{(y)}
\NormalTok{\}}
\end{Highlighting}
\end{Shaded}

To verify the sampling method, we drawn \(n=10000\) samples from \(f\)
and plot its histogram along with a plot of its density function \(f\).

\begin{Shaded}
\begin{Highlighting}[]
\NormalTok{data }\OtherTok{\textless{}{-}} \FunctionTok{rf}\NormalTok{(}\DecValTok{10000}\NormalTok{)}
\NormalTok{ys }\OtherTok{\textless{}{-}} \FunctionTok{seq}\NormalTok{(}\DecValTok{2}\NormalTok{,}\DecValTok{10}\NormalTok{,}\AttributeTok{by=}\NormalTok{.}\DecValTok{1}\NormalTok{)}
\FunctionTok{hist}\NormalTok{(data,}\AttributeTok{freq=}\NormalTok{F,}\AttributeTok{breaks=}\DecValTok{50}\NormalTok{,}\AttributeTok{main=}\StringTok{"f"}\NormalTok{)}
\FunctionTok{lines}\NormalTok{(}\AttributeTok{x=}\NormalTok{ys,}\FunctionTok{f}\NormalTok{(ys),}\AttributeTok{col=}\StringTok{"red"}\NormalTok{) }
\end{Highlighting}
\end{Shaded}

\includegraphics{exam1_files/figure-latex/unnamed-chunk-2-1.pdf}

The histogram seems compatible with the pdf \(f\).

\fbox{%
\begin{minipage}{.8\textwidth}
NOTE TO Dr. Q: On the in-class portion, I forgot to plug the value that I found
that maximizes $h$ back into $h$ to find $c$. Maybe my brain is just getting too old!
\end{minipage}}

\hypertarget{problem-2-5-points}{%
\section{Problem 2 (5 points)}\label{problem-2-5-points}}

\fbox{%
\begin{minipage}{.8\textwidth}
Implement your accept-reject algorithm to get a sample of 10000 from $\rm{Gamma}(2.5, 1)$.
Verify that your method works via a plot of the true normalized density, and a histogram of the
generated values.
\end{minipage}}

This is not the approach you were looking for, but on the test I did
poorly on this section.

Be that as it may, here is essentially my answer on the test, except
that the solution is so slow that I decreased \(c\) so that, at the
extreme tail of the distribution, \(f(y) / c g(y) > 1\).

\begin{Shaded}
\begin{Highlighting}[]
\NormalTok{alpha }\OtherTok{\textless{}{-}} \FloatTok{2.5}
\NormalTok{n }\OtherTok{\textless{}{-}} \DecValTok{10000}
\NormalTok{c }\OtherTok{\textless{}{-}} \DecValTok{10}
\NormalTok{a }\OtherTok{\textless{}{-}} \FunctionTok{floor}\NormalTok{(alpha)}
\NormalTok{b }\OtherTok{\textless{}{-}}\NormalTok{ a }\SpecialCharTok{/}\NormalTok{ alpha}

\NormalTok{ys }\OtherTok{\textless{}{-}} \FunctionTok{vector}\NormalTok{(}\AttributeTok{length=}\NormalTok{n)}
\ControlFlowTok{for}\NormalTok{ (i }\ControlFlowTok{in} \DecValTok{1}\SpecialCharTok{:}\NormalTok{n)}
\NormalTok{\{}
  \ControlFlowTok{repeat}
\NormalTok{  \{}
\NormalTok{    y }\OtherTok{\textless{}{-}} \FunctionTok{rgamma}\NormalTok{(}\AttributeTok{n=}\DecValTok{1}\NormalTok{,}\AttributeTok{shape=}\NormalTok{a,}\AttributeTok{scale=}\NormalTok{b)}
\NormalTok{    u }\OtherTok{\textless{}{-}} \FunctionTok{runif}\NormalTok{(}\DecValTok{1}\NormalTok{)}
\NormalTok{    P }\OtherTok{\textless{}{-}} \FunctionTok{dgamma}\NormalTok{(y,}\AttributeTok{shape=}\NormalTok{alpha,}\AttributeTok{scale=}\DecValTok{1}\NormalTok{) }\SpecialCharTok{/}\NormalTok{ (c}\SpecialCharTok{*}\FunctionTok{dgamma}\NormalTok{(y,}\AttributeTok{shape=}\NormalTok{a,}\AttributeTok{scale=}\NormalTok{b))}
    \ControlFlowTok{if}\NormalTok{ (u }\SpecialCharTok{\textless{}}\NormalTok{ P)}
\NormalTok{    \{}
\NormalTok{      ys[i] }\OtherTok{\textless{}{-}}\NormalTok{ y}
      \ControlFlowTok{break}
\NormalTok{    \}}
\NormalTok{  \}}
\NormalTok{\}}

\FunctionTok{hist}\NormalTok{(ys,}\AttributeTok{freq=}\NormalTok{F,}\AttributeTok{breaks=}\DecValTok{50}\NormalTok{)}
\FunctionTok{lines}\NormalTok{(}\AttributeTok{x=}\FunctionTok{seq}\NormalTok{(.}\DecValTok{01}\NormalTok{,}\DecValTok{10}\NormalTok{,}\AttributeTok{by=}\NormalTok{.}\DecValTok{05}\NormalTok{),}\AttributeTok{y=}\FunctionTok{dgamma}\NormalTok{(}\FunctionTok{seq}\NormalTok{(.}\DecValTok{01}\NormalTok{,}\DecValTok{10}\NormalTok{,}\AttributeTok{by=}\NormalTok{.}\DecValTok{05}\NormalTok{),}\AttributeTok{shape=}\NormalTok{alpha,}\AttributeTok{scale=}\DecValTok{1}\NormalTok{),}\AttributeTok{col=}\StringTok{"red"}\NormalTok{)}
\end{Highlighting}
\end{Shaded}

\includegraphics{exam1_files/figure-latex/unnamed-chunk-3-1.pdf}

\hypertarget{problem-3-30-points}{%
\section{Problem 3 (30 points)}\label{problem-3-30-points}}

\fbox{%
\begin{minipage}{.8\textwidth}
Considers $197$ animals randomly divided into four categories (four phenotypes)
as follows:
$X = (x1, x2, x3, x4)^T$ with cell probabilities $(1/2 + θ/4,(1 − θ)/4,(1 − θ)/4, θ/4)^T$.
We observe $X = (125, 18, 20, 34)^T$.
\end{minipage}}

\hypertarget{part-a-1}{%
\subsection{Part (a)}\label{part-a-1}}

\fbox{%
\begin{minipage}{.8\textwidth}
Hand code your algorithm using Newton-Ralphson to find the maximum likelihood
estimator of $θ$ directly from the observed likelihood. Compare your result with
what you get using the built-in $\operatorname{optim}()$ function in R.
\end{minipage}}

In part 1 of the exam, the log-likelihood given the data was determined
to be \[
  \ell(\theta | \vec{x}) = x_1 \log\left(\frac{1}{2}+\frac{\theta}{4}\right) + x_2 \log\left(\frac{1-\theta}{4}\right) + x_3 \log\left(\frac{1-\theta}{4}\right) + x_4 \log\left(\frac{\theta}{4}\right).
\]

We derived the updating equation for the MLE of \(\theta\) to be \[
  \theta^{(t+1)} = \theta^{(t)} - \frac{d \ell / d \theta}{d^2 \ell / d \theta^2}
\] where \[
  \frac{d \ell}{d \theta} = \frac{x_1}{\theta+2} + \frac{x_2}{\theta-1} + \frac{x_3}{\theta-1} + \frac{x_4}{\theta}
\] and \[
  \frac{d^2 \ell}{d \theta^2} = -\frac{x_1}{(\theta+2)^2} - \frac{x_2}{(\theta-1)^2} - \frac{x_3}{(\theta-1)^2} - \frac{x_4}{\theta^2}.
\]

Instead of taking the time to simplify this expression, we will just
substitute these derivations into the updating equation in the following
R code:

\begin{Shaded}
\begin{Highlighting}[]
\NormalTok{l1 }\OtherTok{\textless{}{-}} \ControlFlowTok{function}\NormalTok{(theta,x)}
\NormalTok{\{}
\NormalTok{  x[}\DecValTok{1}\NormalTok{]}\SpecialCharTok{/}\NormalTok{(theta}\SpecialCharTok{+}\DecValTok{2}\NormalTok{) }\SpecialCharTok{+}\NormalTok{ (x[}\DecValTok{2}\NormalTok{]}\SpecialCharTok{+}\NormalTok{x[}\DecValTok{3}\NormalTok{])}\SpecialCharTok{/}\NormalTok{(theta}\DecValTok{{-}1}\NormalTok{) }\SpecialCharTok{+}\NormalTok{ x[}\DecValTok{4}\NormalTok{]}\SpecialCharTok{/}\NormalTok{theta}
\NormalTok{\}}

\NormalTok{l2 }\OtherTok{\textless{}{-}} \ControlFlowTok{function}\NormalTok{(theta,x)}
\NormalTok{\{}
  \SpecialCharTok{{-}}\NormalTok{x[}\DecValTok{1}\NormalTok{]}\SpecialCharTok{/}\NormalTok{(theta}\SpecialCharTok{+}\DecValTok{2}\NormalTok{)}\SpecialCharTok{\^{}}\DecValTok{2} \SpecialCharTok{{-}}\NormalTok{ (x[}\DecValTok{2}\NormalTok{]}\SpecialCharTok{+}\NormalTok{x[}\DecValTok{3}\NormalTok{])}\SpecialCharTok{/}\NormalTok{(theta}\DecValTok{{-}1}\NormalTok{)}\SpecialCharTok{\^{}}\DecValTok{2} \SpecialCharTok{{-}}\NormalTok{ x[}\DecValTok{4}\NormalTok{]}\SpecialCharTok{/}\NormalTok{theta}\SpecialCharTok{\^{}}\DecValTok{2}
\NormalTok{\}}

\NormalTok{theta\_mle }\OtherTok{\textless{}{-}} \ControlFlowTok{function}\NormalTok{(x, }\AttributeTok{start =} \FloatTok{0.5}\NormalTok{, }\AttributeTok{eps =} \FloatTok{1e{-}6}\NormalTok{)}
\NormalTok{\{}
\NormalTok{  i }\OtherTok{\textless{}{-}} \DecValTok{0}
\NormalTok{  theta0 }\OtherTok{\textless{}{-}}\NormalTok{ start}
\NormalTok{  theta1 }\OtherTok{\textless{}{-}} \ConstantTok{NULL}
  \ControlFlowTok{repeat}
\NormalTok{  \{}
\NormalTok{    theta1 }\OtherTok{\textless{}{-}}\NormalTok{ theta0 }\SpecialCharTok{{-}} \FunctionTok{l1}\NormalTok{(theta0,x) }\SpecialCharTok{/} \FunctionTok{l2}\NormalTok{(theta0,x)}
    \ControlFlowTok{if}\NormalTok{ (}\FunctionTok{abs}\NormalTok{(theta1}\SpecialCharTok{{-}}\NormalTok{theta0) }\SpecialCharTok{\textless{}}\NormalTok{ eps) \{ }\ControlFlowTok{break}\NormalTok{ \}}
\NormalTok{    theta0 }\OtherTok{\textless{}{-}}\NormalTok{ theta1}
\NormalTok{    i }\OtherTok{\textless{}{-}}\NormalTok{ i }\SpecialCharTok{+} \DecValTok{1}
\NormalTok{  \}}
  \FunctionTok{list}\NormalTok{(}\AttributeTok{mle=}\NormalTok{theta1,}\AttributeTok{iterations=}\NormalTok{i)}
\NormalTok{\}}
\end{Highlighting}
\end{Shaded}

We invoke the updating equations on the given data to estimate
\(\theta\) with the following R code:

\begin{Shaded}
\begin{Highlighting}[]
\CommentTok{\# observed data}
\NormalTok{x }\OtherTok{\textless{}{-}}\FunctionTok{c}\NormalTok{(}\DecValTok{125}\NormalTok{, }\DecValTok{18}\NormalTok{, }\DecValTok{20}\NormalTok{, }\DecValTok{34}\NormalTok{)}
\NormalTok{sol }\OtherTok{\textless{}{-}} \FunctionTok{theta\_mle}\NormalTok{(}\AttributeTok{x=}\NormalTok{x, }\AttributeTok{start=}\FloatTok{0.5}\NormalTok{, }\AttributeTok{eps=}\FloatTok{1e{-}6}\NormalTok{)}
\NormalTok{mle }\OtherTok{\textless{}{-}}\NormalTok{ sol}\SpecialCharTok{$}\NormalTok{mle}
\NormalTok{mle\_iterations }\OtherTok{\textless{}{-}}\NormalTok{ sol}\SpecialCharTok{$}\NormalTok{iterations}
\end{Highlighting}
\end{Shaded}

We see that the MLE converges to a solution around \(0.6268215\) after
\(3\). We compare this result with the built-in procedure,
\(\rm{optim}\):

\begin{Shaded}
\begin{Highlighting}[]
\CommentTok{\# the function to maximize, the log{-}likelihood function}
\NormalTok{loglike }\OtherTok{\textless{}{-}} \ControlFlowTok{function}\NormalTok{(theta)}
\NormalTok{\{}
\NormalTok{  x[}\DecValTok{1}\NormalTok{]}\SpecialCharTok{*}\FunctionTok{log}\NormalTok{(}\FloatTok{0.5}\SpecialCharTok{+}\NormalTok{theta}\SpecialCharTok{/}\DecValTok{4}\NormalTok{) }\SpecialCharTok{+}\NormalTok{ x[}\DecValTok{2}\NormalTok{]}\SpecialCharTok{*}\FunctionTok{log}\NormalTok{(}\FloatTok{0.25}\SpecialCharTok{*}\NormalTok{(}\DecValTok{1}\SpecialCharTok{{-}}\NormalTok{theta)) }\SpecialCharTok{+}\NormalTok{ x[}\DecValTok{3}\NormalTok{]}\SpecialCharTok{*}\FunctionTok{log}\NormalTok{(}\FloatTok{0.25}\SpecialCharTok{*}\NormalTok{(}\DecValTok{1}\SpecialCharTok{{-}}\NormalTok{theta)) }\SpecialCharTok{+}\NormalTok{ x[}\DecValTok{4}\NormalTok{]}\SpecialCharTok{*}\FunctionTok{log}\NormalTok{(}\FloatTok{0.25}\SpecialCharTok{*}\NormalTok{theta)}
\NormalTok{\}}

\CommentTok{\# since optim finds the value that minimizes the function, we provide it with}
\CommentTok{\# the negative of the log{-}likelihood.}
\FunctionTok{optim}\NormalTok{(}\FloatTok{0.5}\NormalTok{,}\ControlFlowTok{function}\NormalTok{(theta) \{ }\SpecialCharTok{{-}}\FunctionTok{loglike}\NormalTok{(theta) \},}\AttributeTok{lower=}\FloatTok{0.1}\NormalTok{,}\AttributeTok{upper=}\FloatTok{0.9}\NormalTok{,}\AttributeTok{method=}\StringTok{"L{-}BFGS{-}B"}\NormalTok{)}\SpecialCharTok{$}\NormalTok{par}
\end{Highlighting}
\end{Shaded}

\begin{verbatim}
## [1] 0.626821
\end{verbatim}

The value \(\rm{optim}\) found is approximately the same. It was fussy
with explicitly providing upper and lower bounds. This warrants further
investigation, but for another time.

\hypertarget{part-b-1}{%
\subsection{Part (b)}\label{part-b-1}}

\fbox{%
\begin{minipage}{.8\textwidth}
Implement the E-M algorithm to find the the maximum likelihood estimator of $θ$
with the ``augmented'' data with missing information $Z = y_2$. Compare your
result with part (a).
\end{minipage}}

We split the first category into two, \[
  \vec{y} = (y_1, y_2, y_3, y_4, y_5)^T
\] with cell probabilities \((1/2, θ/4,(1 − θ)/4,(1 − θ)/4, θ/4)^T\).

So, \(y_1\) and \(y_2\) are \emph{unobserved}. The pdf is given by \[
  f(\vec{y}|\theta) \propto \prod_{i=1}^{5} \pi_i(\theta)^{y_i}
\] where \[
  \vec{\pi}(\theta) = (1/2, θ/4,(1 − θ)/4,(1 − θ)/4, θ/4)^T.
\]

Thus, the complete log-likelihood is given by \[
  \log f(\vec{y} | \theta) = \sum_{i=1}^{5} y_i \log \pi_i(\theta)
\] which may be written as \[
  \log f(\vec{y} | \theta) = k(\vec{y}) + y_2 \log(\theta) + y_3 \log(1-\theta) + y_4 \log(1-\theta) + y_5 \log \theta.
\]

\hypertarget{e-step}{%
\subsubsection{E-step}\label{e-step}}

The function we seek to maximize \(Q\) is thus given by \[
  Q(\theta|\theta^{(t)}) = E_{Z|\vec{x},\theta^{(t)}}(\log f(\vec{y}|\theta)),
\] which is given by \[
  Q(\theta|\theta^{(t)})= E_{Z|\vec{x},|\theta^{(t)}}\left(y_2 \log(\theta) + y_3 \log(1-\theta) + y_4 \log(1-\theta) + y_5 \log \theta\right)
\] which by the property of linearity of expectations is equivalent to
\begin{equation}
\begin{split}
Q(\theta|\theta^{(t)})
  = E_{Z|\vec{x},|\theta^{(t)}}(y_2) & \log(\theta) + E_{Z|\vec{x},|\theta^{(t)}}(y_3) \log(1-\theta) +\\
    & E_{Z|\vec{x},|\theta^{(t)}}(y_4) \log(1-\theta) + E_{Z|\vec{x},|\theta^{(t)}}(y_5) \log \theta.
\end{split}
\end{equation}

Note that \(y_3\), \(y_4\), and \(y_5\) are respectively observed to be
\(x_2\),\(x_3\),\(x_4\) and thus, for instance, \[
  E_{Z|\vec{x}}(y_3) = x_2.
\]

The missing data \(Z = y_2\) is distributed \[
  Z | (\vec{x},\theta^{(t)}) \sim \rm{BIN}(x_1, \frac{\theta^{(t)}}{2+\theta^{(t)}}).
\] and thus \[
  E_{Z|\vec{x},|\theta^{(t)}}(y_2) = \frac{x_1 \theta^{(t)}}{2+\theta^{(t)}}.
\]

Putting it all together, we may rewrite \(Q\) as \[
  Q(\theta|\theta^{(t)}) = \frac{x_1 \theta^{(t)}}{2+\theta^{(t)}} \log(\theta) + x_2 \log(1-\theta) + x_3 \log(1-\theta) + x_4 \log \theta.
\]

\hypertarget{m-step}{%
\subsubsection{M-step}\label{m-step}}

In the M-step, we seek \[
 \theta^{(t+1)} = \operatorname{arg\,max}_{\theta} Q(\theta|\theta^{(t)}),
\] which can be found by solving \[
  \left. \frac{d Q(\theta|\theta^{(t)})}{d \theta} \right\vert_{\theta = \theta^{(t+1)}} = 0.
\]

Since the final solution is a little noisy when we leave \(\vec{x}\)
symbolic, we are going to finish the solution with the observed values
of \(\vec{x}\). So, \[
  \frac{d Q}{d \theta} = \frac{125 \theta^{(t)}}{2+\theta^{(t)}} - \frac{38}{1-\theta}
    - \frac{34}{\theta} = 0
\] has the solution, when we replace \(\theta\) with \(\theta^{(t+1)}\),
\[
  \theta^{(t+1)} = \frac{159 \theta^{(t)} + 68}{197 \theta^{(t)} + 144}.
\]

Iterating update equation until some stopping condition that signifies
convergence may then be done. We denote the converged solution of the EM
algorithm by \[
  \hat\theta_{\rm{EM}} = \lim_n \theta^{(t)}.
\]

We implement the EM algorithm in the following R code:

\begin{Shaded}
\begin{Highlighting}[]
\NormalTok{theta0 }\OtherTok{\textless{}{-}}\NormalTok{ .}\DecValTok{5}
\NormalTok{theta1 }\OtherTok{\textless{}{-}} \ConstantTok{NULL}
\NormalTok{i }\OtherTok{\textless{}{-}} \DecValTok{1}
\ControlFlowTok{repeat}
\NormalTok{\{}
\NormalTok{  theta1 }\OtherTok{\textless{}{-}}\NormalTok{ (}\DecValTok{159}\SpecialCharTok{*}\NormalTok{theta0 }\SpecialCharTok{+} \DecValTok{68}\NormalTok{)}\SpecialCharTok{/}\NormalTok{(}\DecValTok{197}\SpecialCharTok{*}\NormalTok{theta0 }\SpecialCharTok{+} \DecValTok{144}\NormalTok{)}
  \FunctionTok{cat}\NormalTok{(}\StringTok{"theta["}\NormalTok{,i,}\StringTok{"] = "}\NormalTok{, theta1,}\StringTok{"}\SpecialCharTok{\textbackslash{}n}\StringTok{"}\NormalTok{)}
  \ControlFlowTok{if}\NormalTok{ (}\FunctionTok{abs}\NormalTok{(theta1}\SpecialCharTok{{-}}\NormalTok{theta0) }\SpecialCharTok{\textless{}} \FloatTok{1e{-}6}\NormalTok{) \{ }\ControlFlowTok{break}\NormalTok{ \}}
\NormalTok{  theta0 }\OtherTok{\textless{}{-}}\NormalTok{ theta1}
\NormalTok{  i }\OtherTok{\textless{}{-}}\NormalTok{ i }\SpecialCharTok{+} \DecValTok{1}
\NormalTok{\}}
\end{Highlighting}
\end{Shaded}

\begin{verbatim}
## theta[ 1 ] =  0.6082474 
## theta[ 2 ] =  0.6243211 
## theta[ 3 ] =  0.6264889 
## theta[ 4 ] =  0.6267773 
## theta[ 5 ] =  0.6268156 
## theta[ 6 ] =  0.6268207 
## theta[ 7 ] =  0.6268214
\end{verbatim}

\begin{Shaded}
\begin{Highlighting}[]
\FunctionTok{print}\NormalTok{(theta1)}
\end{Highlighting}
\end{Shaded}

\begin{verbatim}
## [1] 0.6268214
\end{verbatim}

We see the standard MLE \(\hat\theta_{\rm{mle}}\) and
\(\hat\theta_{\rm{EM}}\) both obtain the same solution up to \(6\)
decimal places when using the same starting value and stopping
condition. However, \(\hat\theta_{\rm{EM}}\) required more iterations
before convergence, which was expected given that the EM algorithm is of
linear order while the MLE using Newton-raphson is of quadratic order.
However, the EM algorithm does have the benefit of a less complex
updating equation.

\hypertarget{part-c}{%
\subsection{Part (c)}\label{part-c}}

\fbox{%
\begin{minipage}{.8\textwidth}
Find the standard error of the MLE using either a numerical calculation of the inverse of the
observed fisher’s information, i.e $[−l”(ˆθ)]^{−1}$, or using Louis’ Method.
\end{minipage}}

We already have the second derivative of the log-likelihood, so we
choose to use the observed Fisher information evaluated at
\(\theta = \theta_{\rm{mle}}\).

\begin{Shaded}
\begin{Highlighting}[]
\NormalTok{obs\_fisher }\OtherTok{\textless{}{-}} \SpecialCharTok{{-}}\FunctionTok{l2}\NormalTok{(mle,x)}
\NormalTok{var\_mle }\OtherTok{\textless{}{-}} \DecValTok{1} \SpecialCharTok{/}\NormalTok{ obs\_fisher}
\FunctionTok{sqrt}\NormalTok{(var\_mle)}
\end{Highlighting}
\end{Shaded}

\begin{verbatim}
## [1] 0.05146735
\end{verbatim}

We see that an estimate of standard error of the MLE is
\(\operatorname{sd}(\hat{\theta}_{\rm{mle}}) = 0.0514673\).

\hypertarget{problem-4-10-points}{%
\section{Problem 4 (10 points )}\label{problem-4-10-points}}

\fbox{Consider the density $f(x) ∝ 3e^{−0.5(x+2)^2} + 7e^{−0.5(x−2)^2}$ (problems 6, Homework 1).}

\hypertarget{part-a-2}{%
\subsection{Part (a)}\label{part-a-2}}

\fbox{%
\begin{minipage}{.8\textwidth}
Compute the exact normalizing constant, both in closed form using $\pi$ (pencil and paper), and
also to 5 decimal places (by evaluating the exact version in R).
\end{minipage}}

It is given that \[
  f(x) \propto \ker(x)
\] i.e., \(f(x) = C \ker(x)\), where \[
  \rm{ker}(x) = 3 e^{-0.5(x+2)^2} + 7 e^{-0.5(x-2)^2}
\] and \(C\) is the normalizing constant that satisfies the equation \[
  \int_{-\infty}^{\infty} \frac{1}{C} \ker(x) dx = 1.
\]

Solving for \(C\) yields the result \[
  C = \int_{-\infty}^{\infty} 3 e^{-0.5(x+2)^2} dx + \int_{-\infty}^{\infty} 7 e^{-0.5(x-2)^2} dx.
\]

Each mode of the bimodal distribution resembles a normal distribution
with \(\sigma = 1\). We may rewrite the above to \[
  C = 3 \sqrt{2 \pi} \int_{-\infty}^{\infty} \frac{e^{-0.5(x+2)^2}}{\sqrt{2 \pi}} dx + 7 \sqrt{2 \pi} \int_{-\infty}^{\infty} \frac{e^{-0.5(x-2)^2}}{\sqrt{2 \pi}} dx
\] which is equivalent to \[
  C = 3 \sqrt{2 \pi} \int_{-\infty}^{\infty} \phi(x+2) dx + 7 \sqrt{2 \pi} \int_{-\infty}^{\infty} \phi(x-2) dx.
\] Since \(\phi\) is a pdf, each integral evaluates to \(1\), and thus
\[
  C = 3 \sqrt{2 \pi} + 7 \sqrt{2 \pi} = 10 \sqrt{2 \pi}
\] is the normalizing constant, which is 25.06628 to \(5\) decimal
places.

Using R's numerical integrator, we get the result:

\begin{Shaded}
\begin{Highlighting}[]
\NormalTok{C }\OtherTok{\textless{}{-}} \DecValTok{10}\SpecialCharTok{*}\FunctionTok{sqrt}\NormalTok{(}\DecValTok{2}\SpecialCharTok{*}\NormalTok{pi)}
\NormalTok{ker }\OtherTok{\textless{}{-}} \ControlFlowTok{function}\NormalTok{(x) \{ }\DecValTok{3}\SpecialCharTok{*}\FunctionTok{exp}\NormalTok{(}\SpecialCharTok{{-}}\FloatTok{0.5}\SpecialCharTok{*}\NormalTok{(x}\SpecialCharTok{+}\DecValTok{2}\NormalTok{)}\SpecialCharTok{\^{}}\DecValTok{2}\NormalTok{) }\SpecialCharTok{+} \DecValTok{7}\SpecialCharTok{*}\FunctionTok{exp}\NormalTok{(}\SpecialCharTok{{-}}\FloatTok{0.5}\SpecialCharTok{*}\NormalTok{(x}\DecValTok{{-}2}\NormalTok{)}\SpecialCharTok{\^{}}\DecValTok{2}\NormalTok{) \}}
\NormalTok{res }\OtherTok{\textless{}{-}} \FunctionTok{integrate}\NormalTok{(ker,}\AttributeTok{lower =} \SpecialCharTok{{-}}\ConstantTok{Inf}\NormalTok{, }\AttributeTok{upper =} \ConstantTok{Inf}\NormalTok{)}
\FunctionTok{print}\NormalTok{(res)}
\end{Highlighting}
\end{Shaded}

\begin{verbatim}
## 25.06628 with absolute error < 4e-05
\end{verbatim}

\hypertarget{part-b-2}{%
\subsection{Part (b)}\label{part-b-2}}

\fbox{%
\begin{minipage}{.8\textwidth}
Approximate this integral with a Simpson’s rule to three decimal places.
What is the effective range and number of subintervals required?
\end{minipage}}

Simpson's rule is implemented by the following R code. Note that we did
not bother to optimize it.

\begin{Shaded}
\begin{Highlighting}[]
\CommentTok{\# simpson : numerical integrator applying simpson\textquotesingle{}s rule to n subintervals}
\CommentTok{\# over the range (a,b).}
\CommentTok{\#}
\CommentTok{\# arguments;}
\CommentTok{\#   f : the function to integrate}
\CommentTok{\#   a : the lower{-}bound}
\CommentTok{\#   b : the upper{-}bound. (a,b) together discrete the range.}
\CommentTok{\#   n : the number of subintervals to partition the range}
\CommentTok{\#}
\CommentTok{\# we evenly partition the range into n subintervals of size (b{-}a)/n}
\CommentTok{\# and apply simson\textquotesingle{}s rule to each subinterval. then, we accumulate these}
\CommentTok{\# values and return the result.}
\NormalTok{simpson }\OtherTok{\textless{}{-}} \ControlFlowTok{function}\NormalTok{(f, a, b, n)}
\NormalTok{  \{}
\NormalTok{    h }\OtherTok{\textless{}{-}}\NormalTok{ (b}\SpecialCharTok{{-}}\NormalTok{a)}\SpecialCharTok{/}\NormalTok{n}
\NormalTok{    s }\OtherTok{\textless{}{-}} \DecValTok{0}
\NormalTok{    x }\OtherTok{\textless{}{-}}\NormalTok{ a}
    \ControlFlowTok{for}\NormalTok{ (i }\ControlFlowTok{in} \DecValTok{1}\SpecialCharTok{:}\NormalTok{(n}\SpecialCharTok{/}\DecValTok{2}\NormalTok{))}
\NormalTok{    \{}
\NormalTok{      s }\OtherTok{\textless{}{-}}\NormalTok{ s }\SpecialCharTok{+} \FunctionTok{f}\NormalTok{(x) }\SpecialCharTok{+} \DecValTok{4} \SpecialCharTok{*} \FunctionTok{f}\NormalTok{(x}\SpecialCharTok{+}\NormalTok{h) }\SpecialCharTok{+} \FunctionTok{f}\NormalTok{(x}\SpecialCharTok{+}\DecValTok{2}\SpecialCharTok{*}\NormalTok{h)}
\NormalTok{      x }\OtherTok{\textless{}{-}}\NormalTok{ x }\SpecialCharTok{+} \DecValTok{2}\SpecialCharTok{*}\NormalTok{h}
\NormalTok{    \}}
\NormalTok{    s}\SpecialCharTok{*}\NormalTok{h}\SpecialCharTok{/}\DecValTok{3}
\NormalTok{\}}
\end{Highlighting}
\end{Shaded}

To test effective range and subintervals, we decided to make a program
that exhaustively searches for the minimum range and number of
subintervals over a discrete set of points. In particular, we search for
a range of the form \((-r,r)\) that is symmetric and an even \(n\). It
is not perfect, but it seems like a reasonable way to estimate these
requirements.

Here is the R code:

\begin{Shaded}
\begin{Highlighting}[]
\NormalTok{R }\OtherTok{\textless{}{-}} \ConstantTok{NULL}
\NormalTok{N }\OtherTok{\textless{}{-}} \ConstantTok{NULL}
\NormalTok{found }\OtherTok{\textless{}{-}}\NormalTok{ F}
\ControlFlowTok{for}\NormalTok{ (r }\ControlFlowTok{in} \DecValTok{1}\SpecialCharTok{:}\DecValTok{200}\NormalTok{)}
\NormalTok{\{}
  \ControlFlowTok{for}\NormalTok{ (n }\ControlFlowTok{in} \DecValTok{1}\SpecialCharTok{:}\DecValTok{100}\NormalTok{)}
\NormalTok{  \{}
\NormalTok{    res }\OtherTok{\textless{}{-}} \FunctionTok{simpson}\NormalTok{(ker,}\SpecialCharTok{{-}}\NormalTok{r}\SpecialCharTok{/}\DecValTok{10}\NormalTok{,r}\SpecialCharTok{/}\DecValTok{10}\NormalTok{,}\DecValTok{2}\SpecialCharTok{*}\NormalTok{n)}
    \ControlFlowTok{if}\NormalTok{ (}\FunctionTok{abs}\NormalTok{(res }\SpecialCharTok{{-}}\NormalTok{ C) }\SpecialCharTok{\textless{}} \FloatTok{0.001}\NormalTok{)}
\NormalTok{    \{}
\NormalTok{      R }\OtherTok{\textless{}{-}}\NormalTok{ r}\SpecialCharTok{/}\DecValTok{10}
\NormalTok{      N }\OtherTok{\textless{}{-}} \DecValTok{2}\SpecialCharTok{*}\NormalTok{n}
      \FunctionTok{cat}\NormalTok{(}\StringTok{"r = +{-}"}\NormalTok{, R, }\StringTok{", n = "}\NormalTok{, N, }\StringTok{", C = "}\NormalTok{, res, }\StringTok{"}\SpecialCharTok{\textbackslash{}n}\StringTok{"}\NormalTok{)}
\NormalTok{      found }\OtherTok{\textless{}{-}}\NormalTok{ T}
      \ControlFlowTok{break}
\NormalTok{    \}}
\NormalTok{  \}}
  \ControlFlowTok{if}\NormalTok{ (found) \{ }\ControlFlowTok{break}\NormalTok{ \}}
\NormalTok{\}}
\end{Highlighting}
\end{Shaded}

\begin{verbatim}
## r = +- 5.9 , n =  16 , C =  25.06603
\end{verbatim}

The minimum range \((-r,r)\), \(r > 0\), when divided into \(n\)
subintervals, that is the same as the true value \(10 \sqrt{2\pi}\) to
\(3\) decimal places is given by \(r = 5.9\) and \(n = 16\).

\end{document}
